En este capítulo se describe el desarrollo del sistema siguiendo las fases de
la metodología Mobile-D presentadas en la sección~\ref{sec:mt-mobiled}. Para
cada fase se detallan las actividades realizadas y los artefactos obtenidos.
La fase de producción se organizó en seis iteraciones, una por módulo
funcional, y cada iteración se documenta con su planificación, el trabajo de
diseño e implementación y la liberación con sus criterios de aceptación.

El producto resultante se identifica en este documento como \nombreApp{}, el
nombre de la aplicación registrado en su configuración de compilación. El
código fuente se gestiona en un repositorio Git alojado en GitHub, y los
fragmentos de código que se presentan corresponden a la versión integrada en
la rama \archivo{develop}.

La Tabla~\ref{tab:resumen-fases} resume la correspondencia entre las fases de
Mobile-D, las actividades principales y los artefactos producidos.

\begin{table}[H]
\centering
\caption{Resumen de las fases de Mobile-D aplicadas al proyecto}
\label{tab:resumen-fases}
\begin{tabularx}{\textwidth}{|L{2.6cm}|Y|Y|}
\hline
\rowcolor{encabezado} \textbf{Fase} & \textbf{Actividades principales} & \textbf{Artefactos} \\
\hline
Exploración & Identificación de interesados, definición de alcance,
requerimientos, selección tecnológica y planificación. & Requerimientos
funcionales y no funcionales, casos de uso, plan de iteraciones. \\ \hline
Inicialización & Configuración del entorno, repositorio, proyecto Firebase,
arquitectura, modelo de datos y día de prueba. & Arquitectura por capas,
modelo de datos de Firestore, reglas de seguridad, estructura del
proyecto. \\ \hline
Producción & Seis iteraciones: autenticación, guía de máquinas, escáner,
rutinas, nutrición y progreso. & Módulos funcionales integrados por
\emph{pull request}. \\ \hline
Estabilización & Integración, refactorización por capas, validación y
verificación estática. & Código uniforme y verificado. \\ \hline
Pruebas del sistema & Pruebas funcionales, de seguridad y de aceptación por
objetivo específico. & Matriz de casos de prueba y resultados. \\
\hline
\end{tabularx}
\end{table}
